\documentclass[12pt]{article}
 \usepackage[margin=1in]{geometry} 
\usepackage{amsmath,amsthm,amssymb,amsfonts}
\usepackage{centernot}
\usepackage[table]{xcolor}
\usepackage{algpseudocode}
\usepackage{algorithm}
\usepackage{graphicx}
\usepackage{amsmath}
\usepackage{graphicx}
\graphicspath{ {./Surveys/} }
\floatname{algorithm}{}
\author{Joshua Stephenson-Losey,
Keefer Sands,
Taylor Koth}
 
\newcommand{\N}{\mathbb{N}}
\newcommand{\Z}{\mathbb{Z}}
\renewcommand*{\proofname}{Solution}


\newenvironment{problem}[2][Problem]
{\begin{trivlist}
\item[\hskip \labelsep {\bfseries #1}\hskip \labelsep {\bfseries #2.}]}{\end{trivlist}}


\makeatletter
\renewcommand{\fnum@algorithm}{\fname@algorithm}
\makeatother

\begin{document}

 
\title{CSCI 432\\Homework 6}
\date{}
\maketitle

Assigned 11/26/2018, due by end of class (4:00 pm) on 12/03/2018. Marked questions (*) are graded for correctness. The remaining will be graded for effort. Please see the course website for details about expected effort.

You must follow the collaboration policy detailed on the course website. Please type solutions in an appropriate editor (\LaTeX, Word) so that I can review equations and proofs efficiently. 

\begin{problem}{1*}
Given a graph $G=(V,E)$, we want to color each node with one of three colors. An edge is \textit{satisfied} if the colors for each of its nodes are different. The goal is to maximize the number of satisfied edges.

\begin{itemize}
\item Design an algorithm that randomly assigns each node a color.

-For each node create a random integer between 1 and 3 (inclusive) and assign the node to a color using this number (ex. 1 = blue.... ect.)

\item Show that the probability a given edge is satisfied is $\frac{2}{3}$.

-Looking at an edge with two nodes and three possible colors for each node gives 9 total combinations of colors, only three of which are invalid. With total probability of one this makes the probability of any given edge being satisfied $\frac{2}{3}$.

\item Use the following fact to show that the expected number of edges satisfied by our algorithm is at least $\frac{2}{3}$OPT. 

Fact: For random variables $X$ and $Y$, $E[X+Y] = E[X]+E[Y]$.

- Let X represent a random single edge-vertex group that equals one if the edge is satisfied and zero otherwise

- The expected value of X is:

$\qquad\frac{2}{3} *$
$1_{(probability\_of\_being\_satisfied)} + \frac{1}{3} *$
$0_{(probability\_of\_not\_being\_satisfied)}$

- The expected value of each edge is then equal to $\frac{2}{3}$

- The expected value of n edges is then $\frac{2}{3}$ * n

- Since the summation of expected values is equal to the expected value of the summation the expected value of all edges able to be satisfied in graph G is $\frac{2}{3}$ * n

- Since OPT will chose the maximum number of edges able to be satisfied in graph G OPT = n

- $\therefore ALG \ge \frac{2}{3} OPT$
\end{itemize}
\end{problem}

\begin{problem}{2}
Suppose you have a random number generator that outputs $1$ with probability $0 < p \le 1$ and $0$ otherwise. Using this as a subroutine, make an algorithm that outputs $1$ with probability $\frac{1}{2}$ and $0$ otherwise.
\end{problem}

\begin{problem}{3*}
Evaluate the course by using the CS survey system: 

https://www.cs.montana.edu/survey/

I will not see what you write until after final grades are submitted. Even then, it is annonymous, so please be honest. I can't improve unless you tell me what is wrong. Put in a screen shot showing it is completed for credit.

Joshua Stephenson-Losey:

\includegraphics[scale = .33]{Surveys/SL_432.png}

\includegraphics[scale = .33]{Surveys/SL_432_yaw.png}

Keefer Sands:

Taylor Koth:


\end{problem}

\begin{problem}{4}
Show up and participate in the workshop session on 11/30.
\end{problem}
 


\end{document}